% Преамбула

\usepackage[T2A]{fontenc}

% В книге говориться, что можно использовать koi8-r (для Linux)
% и cp1251 для Windows
% Я выбрал utf8. .tex файлы сохраняю в этой кодировке. Работает на Linux и Windows.
\usepackage[utf8]{inputenc}

% Пакет babel адаптирует LATEX к работе с русским языком
\usepackage[russian]{babel}

% packages ====================================================================

% [showframe] - useful option
\usepackage[margin=2cm]{geometry}

% Отступ (красная строка) для первого абзаца главы или параграфа
\usepackage{indentfirst}

\usepackage{setspace}
\usepackage[indent]{parskip}

\usepackage{epigraph}

% Оглавление и ссылки как гиперссылки
% Это для удобства чтения в электронном виде
\usepackage[unicode]{hyperref}

% Выравнивание подписей по левому краю
\usepackage{caption}
\captionsetup{justification=raggedright, singlelinecheck=false}

% Для таблиц
\usepackage{makecell}

% Для вставки заранее подготовленных рисунков
\usepackage{graphicx}

% Листинги программ
\usepackage{verbatim}
\usepackage{moreverb}
\usepackage{listings}

\lstdefinestyle{texstyle}
{
    breaklines=true,
    frame=single,
    inputencoding=utf8,
    language=TeX,
    numbers=left,
    % для нормального отображения кириллицы в комментариях
    extendedchars=\true,
    keepspaces=true
}

\lstset{style=texstyle}

% Математические формулы
\usepackage{amsmath}
\usepackage{amsfonts}

\usepackage{datetime}

% Для устранения ворнингов про \circle, \oval...
\usepackage{pict2e}

% layout ====================================================================

% Предотвращение залезания строк на поля, даже если для этого требуется заполнить строку недопустимо длинными пробелами
\sloppy

% Установка межстрочного интервала
% \onehalfspacing

% Установка интервала между абзацами
% \setlength{\parskip}{1em}

% macros ====================================================================
% Определение (sbdef - serge beloff definition)
\newcommand{\sbdef}[1]{\textbf{#1}}

% 0.98 - to fix compiler warning Overfull \hbox (6.79999pt too wide) in paragraph
\newcommand{\deff}[1]{\noindent\fbox{\parbox[l]{0.98\textwidth}{\textbf{#1}}}}

